% Complete full proof; generated without omitting sections.
\section{Whole-expression centers in eleven unrestricted integers}

\subsection{Result and boundary}

Four fixed ordinary integer polynomials have exact (degree, live binary-operation)
pairs (316,133), (308,140), (304,144), (300,148). Each has eleven unrestricted
integer variables and the identical six actual coefficient numerals of round17.
They improve the degree side of the inspected eleven-integer frontier, not its
minimum-operation side. In particular the separate incoming Semantic Signs
(924,107)/(840,108) and Odd Helpers (648,115) endpoints remain relevant
low-operation tradeoffs. No global minimum or proof of consistency is claimed.

The arithmetic change is a fixed even radical weight and a polynomial shear of
one coordinate. The complete zero-existence equivalence with the round17
linear-cofactor source is proved below. Its authenticated ordinary
native/compiler/refutation calibration is retained, not replaced by a new
proof-assistant or object-language ZFC certificate. Composing that theorem,
ZFC proves that universal nonvanishing of each displayed polynomial over Z\textasciicircum{}11
is equivalent to syntactic Con(ZFC). This proves neither side true.

\subsection{1. Complete fixed polynomials}

The only variables are j,C,z,h,r,v,w,T,c,n,s in Z. The six ports
beta,K0,MC,MF,cone\_A,cone\_B are fixed numerals, not additional variables.
Their complete compressed little-endian bytes, SHA256 bindings and independent
finite reconstruction recipe are in coefficients/. In particular
beta=2\textasciicircum{}1064921-1, cone\_A=2\textasciicircum{}1113, cone\_B=2\textasciicircum{}666104; MF is already shifted.
Define these ordinary polynomial abbreviations:

\begingroup\small
\begin{verbatim}
J=6j+1, q=beta J+1, W=12w+4, X=qW, Y=v,
E=XY, a=E+Y, A=a+2, H=4a+3, Delta=a^2+H=A^2-1,
F=(K0+W)C-z(q-1),
R=[q(q-F)-C](q^2-1)+(MC+q MF)J,
k=R+1+hE, eta=c-Yk,
Derr=(cone_A F-q)^2+(cone_B C-q)^2,
B0=k^2(q^2-Derr)-4q^2 eta^2,
Budget=Y^2 B0-k^2 q^8, D=Budget-4,
ell=XY^2, Dt=1+ell(ell+1)k^2,
Dg=1+D Dt s^2, P=Dt Dg,
Dm=1+Delta c^2, B=r A Dm, U=r c^2 Dm,
I2=c^4 B=A c^2 U,
Df=1+Delta I2, Q=Delta^2 I2, delta=Q-1,
V=Df(c-R+4A U T)-c, Cy=QV^2-1,
K=X(3c-2X)+2, M=2q^3, Dc=HM,
S=srH M^5,
zc=Dc T+M K+H(v-q^3),
Ac=(zc^2+3Dc^2)(zc^3-19Dc^3).
\end{verbatim}
\endgroup

Choose one row of the following table; L is a fixed numeral.

\begingroup\small
\begin{verbatim}
degree  L   Psi
316     8   [(8Delta+4)U-r c^4]V
308     8   [(8Delta+4)U-r c^4]V+2U^3 T(c^2-2)
304     8   [(8Delta+4)U-r c^4]V+U^2 T[2U(c^2-2)+16A^3]
300    32   [(32Delta+16)U-4r c^4]V
            +U^2 T[8U(c^2-2)+64A^3+3r c^6].
\end{verbatim}
\endgroup

In all cases the entire polynomial is

\begingroup\small
\begin{verbatim}
Z=S(n+Psi)-Ac,
pa=delta Z^2, pb=delta S^2 P, pc=L^2 S^2 B Df Cy,
N=(pa+pb-pc)^2-4pa pb.                         (1)
\end{verbatim}
\endgroup

All displayed names except the eleven declared variables are abbreviations or
fixed numerals. There are no radical, division, variable-exponent, or implicit
existential gates. The literal circuits share the identities

\begingroup\small
\begin{verbatim}
P=Dt+(Budget-4)(Dt s)^2,
Budget=(Yk)^2(q^2-Derr)-q^2[4(Y eta)^2+(kq^3)^2],
A^2=Delta+1.
\end{verbatim}
\endgroup

\subsection{2. The even-weight integral-center lemma}

We prove a slightly stronger statement: for every fixed nonzero even integer L
and every integer polynomial Psi in the ten variables other than n, (1) has the
same projection of its integer zero set onto those ten variables as the
round17 equation. The latter uses L=2 and Psi=(2Delta+1)UV.

The unconditional arithmetic facts used here are proved in the retained
source/LINEAR\_COFACTOR\_ROUND17\_PROOF.md, Sections 2--4, and use no conclusion
from the present norm. They are recalled explicitly to identify all zero and
sign cases. The authentic affine restrictions imply q even, q=2 mod3,
|q|>=16, X=0 mod8, X=2 mod3 and |X|>=4|q|. Thus A is nonzero,
Delta=A\textasciicircum{}2-1>1, Delta=0 mod3, H is odd/nonzero, Dm>=1 and Dt>=1.
The exact modulo-sixteen budget calculation gives D=8 or12 mod16, hence
D is nonzero and nonsquare, P=1 mod4 and P is nonzero. Also Df and delta
are nonzero, Q is even, and delta is odd.

The rational quintic pi(t)=(t\textasciicircum{}2+3)(t\textasciicircum{}3-19) has no rational root. Since
Dc is nonzero, Ac=Dc\textasciicircum{}5 pi(zc/Dc) is nonzero. If S=0, (1) reduces to
(delta Ac\textasciicircum{}2)\textasciicircum{}2, which is nonzero. Therefore S,s,r and B are nonzero
at every new zero.

Put Da=delta Cy and Cstar=B Df Da. The retained elementary sign proof gives
Da>0. For completeness: if c=0, Q=0 and Da=1; if c!=0 and B<0,
both delta and Cy are negative; if c!=0 and B>0, Df>|c|+1, and
V=Df b-c is nonzero for any integer b, so delta,Cy>0. Moreover B Df>0
whenever c!=0. Thus Cstar is nonzero; its only possible negative case is
c=0,B<0.

The four-factor norm identity for (1), divided by nonzero S and multiplied
by delta, gives a rational signed sum

\begingroup\small
\begin{verbatim}
plus or minus L sqrt(Cstar) plus or minus delta sqrt(P).
\end{verbatim}
\endgroup

If exactly one radicand is negative, its nonzero imaginary term cannot cancel.
If both are negative, then c=0, delta=-1 and Cstar=B. Cancellation would
force L\textasciicircum{}2 B=P, impossible modulo4 because L is even and P=1 mod4.
Therefore Cstar and P are positive.

Write each as a square times a positive squarefree integer. Distinct
squarefree radicals are linearly independent over Q. The only possible
nonrational cancellation would require the same nontrivial squarefree part.
But the coefficient of its first radical is even, while that of the second
is odd: L is even, delta is odd and P is odd. Cancellation is impossible.
Consequently both Cstar and P are positive integer squares.

The exact pairwise-coprimality facts of round17 are unchanged:

\begingroup\small
\begin{verbatim}
gcd(B,Df)=gcd(B,Da)=gcd(Df,Da)=1,
Da=1 mod B, Da=A^2 Dm mod Df,
Df=1 mod A, Df=1 mod Dm, Df=1 mod3.
\end{verbatim}
\endgroup

Indeed Df=1 mod B, Q=0 mod B; modulo Df, Q=-Delta,
delta=-A\textasciicircum{}2,V=-c,Cy=-Dm. Since B=r A Dm these give the asserted
coprimalities. Thus |B|,|Df|,Da are squares. As Da>0, B and Df have
one common sign, and Df=1 mod3 excludes a negative square. We obtain

\begingroup\small
\begin{verbatim}
B=u^2, Df=f^2, Da=e^2
\end{verbatim}
\endgroup

with positive u,f,e. Similarly gcd(Dt,Dg)=1 and Dt>=1 yield
Dt=tau\textasciicircum{}2,Dg=g\textasciicircum{}2 with positive tau,g. In particular Budget>=12.
It follows from Y\textasciicircum{}2 B0>k\textasciicircum{}2 q\textasciicircum{}8 that k,Y are nonzero, B0>0 and
|Y|>|q|\textasciicircum{}3; the centered strip then excludes c=0.

Define proof-only I=c\textasciicircum{}2u and Told=uT. Then

\begingroup\small
\begin{verbatim}
f^2-Delta I^2=1, Q=(Delta I)^2,
V=f^2(c-R+4I Told)-c.
\end{verbatim}
\endgroup

The rational y=e/delta satisfies QV\textasciicircum{}2-(Q-1)y\textasciicircum{}2=1. The elementary
adjacent-coefficient descent makes y integral: if b is its reduced
denominator, b\textasciicircum{}2 divides Q-1. After signs, for |V|>1 the integral
unimodular transformation

\begingroup\small
\begin{verbatim}
V'=(2Q-1)V-2(Q-1)y,
y'=(2Q-1)y-2QV
\end{verbatim}
\endgroup

preserves that denominator and the norm, with 0<|V'|<|V|. The first
coordinate is integral because (Q-1)y is integral. Descent reaches
|V|=1, where |y|=1, forcing b=1. This also directly covers the formal
c=0 endpoint, where y=-1, although it is already excluded here.

The two normalized roots in the norm are therefore integers Lufy and
tau g. Hence Z/S is integral. Since Psi is integral, S divides Ac.
All these conclusions hold for arbitrary nonzero even L; no size bound
on Psi and no sign assumption on it was used.

Now retain the ten coordinates of this new zero and put

\begingroup\small
\begin{verbatim}
n_old=Ac/S+2ufy+tau g-(2Delta+1)UV.
\end{verbatim}
\endgroup

This is an integer and makes one old conjugate factor vanish, so it is an
actual round17 zero. Conversely, an old zero gives the same integral
roots and divisibility by the proven round17 theorem. Retain its ten
coordinates and set

\begingroup\small
\begin{verbatim}
n_new=Ac/S+Lufy+tau g-Psi.
\end{verbatim}
\endgroup

One new factor vanishes. These explicit maps prove equality of the
projected zero sets. They are not claimed to be mutually inverse on n,
since a norm may admit several sign branches. For a fixed L, different
centers separately are literal bijective shears of n.

\subsection{3. Calibration and fixed-witness completeness}

The previous section reduces both directions to an actual old zero, not merely
to a list of necessary Pell properties. Thus all outer scales, signs, rank
hypotheses and syntactic calibration are exactly those of the authenticated
old theorem. Its chain, retained verbatim in source/, is as follows.
Carrier divisibility first yields M|H(v-q\textasciicircum{}3), even Y, and
H|q\textasciicircum{}15 Kq\textasciicircum{}5 for Kq=X(X+2ac)-1. The strong/auxiliary pair above and
c\textasciicircum{}2|I recover the main Pell root by absolute folding. Together with
Dt=tau\textasciicircum{}2, the centered strip, |Y|>|q|\textasciicircum{}3 and the authentic packing bounds
|q|(|q|\textasciicircum{}2-1)<R<|q|\textasciicircum{}4-2, signed rank gives c=sign(c) psi\_A(R),
R=2t+1 and |H|>|X|(|X|-3)\textasciicircum{}t. The powered projection, using the
actual enormous base, forces sign(c)X=2\textasciicircum{}R. Modulo3 then fixes c>0,
X=2\textasciicircum{}R, and the old strip fixes k,Y>0. Exact H|K restores
Y=(2v\_old+1)q\textasciicircum{}3. The uniform half-binomial theorem restores the
positive native interval, and fresh odd strong witnesses restore an
actual signed Q41 source zero. Only that actual restored zero is passed
to the retained source34/native/compiler/finite-refutation theorem.

Conversely the canonical source family gives its margin Budget>64.
After choosing its guard pair s,g, the old proof chooses ONE strong pair
I,f with 4s A d c\textasciicircum{}2|I, where Dm=d\textasciicircum{}2, sets
u=I/c\textasciicircum{}2, tstrong=I/(A d c\textasciicircum{}2), r=A tstrong\textasciicircum{}2, and keeps I,f fixed.
The supported-modulus theorem interpolates Told modulo u s r with
that same pair. An intersective root of pi modulo sr gives integral
Ac/S. The old proof never reselects the strong pair after a residue
request. Our final change only replaces its integral n by n\_new above;
therefore it introduces no new interpolation or choice-order obligation.

The inherited ordinary proof boundary is stated, not promoted: these
constructions prove the calibrated reduction, not Con(ZFC), not its
negation, and not an end-to-end new formal certificate.

\subsection{4. Exact degree proof}

All degree statements are for the whole polynomial at the authentic fixed
coefficient tuple, not for generic ports, syntactic circuit bounds or numeric
line samples. The free-scale degrees inherited directly from the displayed
registers are

\begingroup\small
\begin{verbatim}
a,A:3; Delta:6; R,k:5; Budget:18; Dt:18; P:56;
Dm:8; B:12; U:11; Df:22; Q,delta:28; V:37;
Cy:102; S:20; Ac:35; Psi:54; Z:74;
pa:176; pb:124; 4pa pb:300.
\end{verbatim}
\endgroup

Here k's highest form is R's: hXY has degree4, strictly below5.
Put W0=Sn-Ac, so deg W0=35. In a free polynomial ring with independent
A,c,r,T,R, use weights (3,1,1,1,5) and the exact definitions
Delta=A\textasciicircum{}2-1, Dm=1+Delta c\textasciicircum{}2, B=r A Dm, U=r c\textasciicircum{}2 Dm,
Df=1+Delta c\textasciicircum{}4B, V=Df(c-R+4A U T)-c. Define

\begingroup\small
\begin{verbatim}
Dcenter=Psi^2-L^2 B Df V^2,
Ecenter=delta Psi^2-L^2 B Df Cy
       =delta Dcenter-L^2 B Df(V^2-1).             (2)
\end{verbatim}
\endgroup

The following exact weighted leading-form identities hold; every omitted
monomial has strictly smaller displayed weight.

\begingroup\small
\begin{verbatim}
center  deg Dcenter  highest Dcenter                 deg Ecenter
316     90          -128 A^18 c^28 r^6 T^2           118
308     86          -1024 A^19 c^22 r^5 T^2          114
304     84          -48 A^16 c^28 r^6 T^2            112
300     82           4096 A^16 c^26 r^6 T^2          110.
\end{verbatim}
\endgroup

These are finite exact identities, not numerical conjectures. Here is a compact
route to expansion. Put lambda=(8Delta+4)U-r c\textasciicircum{}4, Psi0=lambda V,
h=2U\textasciicircum{}3T(c\textasciicircum{}2-2), b=16A\textasciicircum{}3U\textasciicircum{}2T and e=3r c\textasciicircum{}6U\textasciicircum{}2T. Then

\begingroup\small
\begin{verbatim}
D316=[8Ur(2c^2-c^4)+r^2c^8-64B]V^2,
D308=D316+2Psi0 h+h^2,
D304=D308+2(Psi0+h)b+b^2,
D300=16D304+8(Psi0+h+b)e+e^2.
\end{verbatim}
\endgroup

These follow just by squaring the four displayed centers and the identity
B Df=B+Delta A\textasciicircum{}2 U\textasciicircum{}2. Expanding U=r[c\textasciicircum{}4A\textasciicircum{}2+c\textasciicircum{}2-c\textasciicircum{}4] and
V=(1+(A\textasciicircum{}2-1)c\textasciicircum{}4 rA[1+(A\textasciicircum{}2-1)c\textasciicircum{}2])
  *(c-R+4A r c\textasciicircum{}2[1+(A\textasciicircum{}2-1)c\textasciicircum{}2]T)-c
in these four finite expressions gives the table. check.py independently
expands them in a sparse free ring; CHECK.json supplies every coefficient of
the exact Ecenter remainders. The auxiliary correction in (2) has weight108,
strictly below each leading delta Dcenter. Multiplication by delta's unique
highest form r A\textasciicircum{}7 c\textasciicircum{}6 gives the respective Ecenter leaders

\begingroup\small
\begin{verbatim}
-128 A^25 c^34 r^7 T^2,
-1024 A^26 c^28 r^6 T^2,
-48 A^23 c^34 r^7 T^2,
4096 A^23 c^32 r^7 T^2.
\end{verbatim}
\endgroup

Now use the exact norm cut

\begingroup\small
\begin{verbatim}
pa+pb-pc=delta W0^2+2delta S Psi W0+delta S^2 P
            +S^2 Ecenter.                         (3)
\end{verbatim}
\endgroup

The first three terms have degrees at most98,137,124. The last has respective
exact degrees158,154,152,150 with a unique highest monomial after the actual
substitution A=a+2. Therefore the squared inner terms have degrees
316,308,304,300. In the first three cases the correction4pa pb, degree300,
is smaller. In the last it ties, but its entire highest form has T-degree2,
whereas the inner-square highest monomial has T-degree4. Thus that monomial
cannot cancel and degree300 is exact.

For explicit coefficients put q0=6 beta j, a0=12q0 wv and note
S0=128sr a0 q0\textasciicircum{}15. The certified highest monomials are

\begingroup\small
\begin{verbatim}
degree  exponents (j,w,v,s,r,c,T)       positive coefficient
316     (114,54,54,4,18,68,4)           2^264 3^168 beta^114
308     (116,56,56,4,16,56,4)           2^276 3^172 beta^116
304     (110,50,50,4,18,68,4)           2^246 3^162 beta^110
300     (110,50,50,4,18,64,4)           2^262 3^160 beta^110.
\end{verbatim}
\endgroup

Each exponent sum is its stated degree. Each is nonzero for the actual
beta. In the degree300 case the T\textasciicircum{}4 highest summand is
2\textasciicircum{}52 s\textasciicircum{}4 r\textasciicircum{}18 c\textasciicircum{}64 T\textasciicircum{}4 a0\textasciicircum{}50 q0\textasciicircum{}60. This proof makes no global lower-bound
claim: a degree300 obstruction for the particular chosen norm/guard family
would not exclude a different construction.

\subsection{5. Literal counts and checks}

The live arithmetic counts are, in increasing degree order:

\begingroup\small
\begin{verbatim}
300: 31 add,22 subtract,95 multiply =148;
304: 30 add,22 subtract,92 multiply =144;
308: 28 add,22 subtract,90 multiply =140;
316: 27 add,21 subtract,85 multiply =133.
\end{verbatim}
\endgroup

Every fixed scalar multiplication and every square costs one binary gate;
reuse is free. All eleven variables and six fixed coefficient ports are live.
The operation-count convention is identical to the inspected predecessors.

build.py exports the actual division-free DAGs. check.py does not import the
producer or its DAG class. It independently binds every named register and
all live gates in free polynomial rings; checks the universal four-conjugate
norm, exact center remainders and degree cuts; propagates attained actual
leading coefficients through every node; authenticates all six decompressed
coefficient payloads and reconstructs their complete finite recipe; checks
all copied source pins and liveness/counts; and rejects altered constants,
variables, degrees and selected semantic gates. Its normal and python -O
runs use explicit runtime checks, not removable assert statements.

The checker is a finite algebraic/regression certificate. The unbounded Pell,
interpolation, intersectivity and native/refutation theorems remain exactly
the stated ordinary mathematical inputs. Source proof text, licenses and
actual coefficient payloads are retained for portable replay.
